\documentclass[11pt]{article}

\usepackage[a4paper,margin=1in]{geometry}
\usepackage{amsmath,amssymb}
\usepackage{graphicx}
\usepackage{booktabs}
\usepackage{caption}
\usepackage{hyperref}

\newcommand{\dd}{\,\mathrm{d}}
\newcommand{\bbE}{\mathbf{E}}
\newcommand{\bbH}{\mathbf{H}}
\newcommand{\bbP}{\mathbf{P}}
\newcommand{\bbL}{\mathbf{L}}
\newcommand{\bbx}{\mathbf{x}}
\newcommand{\calH}{\mathcal{H}}

\begin{document}

\subsection{Conserved quantities}
\label{subsec:results_conserved_quantities}

\paragraph{Claim.}
The three-dimensional Poisson-bracket DG discretization preserves the
Hamiltonian character of the lossless Maxwell problem in the sense relevant for
long-time computation: the electromagnetic energy remains non-drifting over many
periods, the discrete divergence constraints remain at roundoff level, and the
errors in the additional integral diagnostics decrease systematically under
$p$- and $h$-refinement.  The quantities monitored in this experiment are global
integrals, not volume averages,
\begin{align}
  \calH_h(t)
    &= \frac12 \int_{\Omega}
       \left(\varepsilon |\bbE_h|^2 + \mu |\bbH_h|^2 \right)\dd \bbx, \\
  Q_{E,h}(t)
    &= \int_{\Omega} \nabla\cdot(\varepsilon \bbE_h)\dd \bbx,
  &
  Q_{H,h}(t)
    &= \int_{\Omega} \nabla\cdot(\mu \bbH_h)\dd \bbx, \\
  \bbP_h(t)
    &= \int_{\Omega} \bbE_h\times \bbH_h \dd \bbx,
  &
  \bbL_h(t)
    &= \int_{\Omega} \bbx\times(\bbE_h\times \bbH_h)\dd \bbx, \\
  \chi_h(t)
    &= \frac12\int_{\Omega}
       \left(\varepsilon \bbE_h\cdot\nabla\times\bbE_h
       + \mu \bbH_h\cdot\nabla\times\bbH_h\right)\dd \bbx .
\end{align}
For scalar diagnostics the plotted error is the absolute difference from the
analytical value.  For momentum diagnostics the plotted error is the Euclidean
norm of the difference between the three integrated components, for example
$\|\bbP-\bbP_h\|_2$ and $\|\bbL-\bbL_h\|_2$.

\paragraph{Experiment.}
The test is a periodic travelling wave in the cuboid
\[
  \Omega = [0,2]\times[0,1]\times[0,1],
\]
with periodic identifications on all opposite faces.  The material parameters
are constant, $\varepsilon=\mu=1$, and the analytical fields are
\begin{align}
  \bbE(\bbx,t) &= \left(0,0,-\sin(kx-\omega t)\right), \\
  \bbH(\bbx,t) &= \left(0,\sin(kx-\omega t),0\right),
\end{align}
with $k=\omega=2\pi$.  Hence the exact optical chirality and both charges are
zero.  Since the wave propagates in the positive $x$ direction, the exact
integrated linear momentum is $\bbP=(1,0,0)$; with the coordinate origin at the
lower corner of the cuboid, the exact integrated angular momentum is
$\bbL=(0,1/2,-1/2)$.

All runs use the distributed periodic Poisson-bracket Maxwell driver with
centered numerical fluxes, a magnetic-first explicit symplectic partitioned
Runge--Kutta rule, CFL number $0.8$, and the Jaskowiec--Sukumar tetrahedral
cubature rule used by the diagnostics.  The $p$-refinement study uses the
structured mesh
\texttt{examples/meshes/periodic\_box\_structured\_2x1x1.vtk}, final time
$T=1000$, and polynomial orders $p=1,2,3,4$ with ESPRK order $p+1$.  The
corresponding run script is
\texttt{examples/run\_periodic\_cuboid\_centered\_flux\_p\_convergence.sh}.
The stored metadata for the visible $p=1$ and $p=2$ runs report one MPI rank,
768 tetrahedra, $h_{\min}=4.4960164\times 10^{-2}$, and time steps
$3.5967342\times10^{-2}$ and $1.7983994\times10^{-2}$, respectively.

The $h$-refinement study fixes $p=1$ and ESPRK order $2$, runs on four MPI
ranks, and refines the structured periodic mesh by factors of two.  The
published panels use the three complete final-time levels shown in
Table~\ref{tab:conserved_h_runs}.  A fourth level, $h_0/8$, is present in the
run directory but its metadata mark it as incomplete; it is therefore not used
for the claims below.

\begin{table}[htbp]
  \centering
  \caption{Reproducibility data for the $h$-refinement conserved-quantity
  experiment.  The data are taken from the run metadata in
  \texttt{output/periodic\_cuboid\_cf\_structured\_p1\_h*}.}
  \label{tab:conserved_h_runs}
  \resizebox{\textwidth}{!}{%
  \begin{tabular}{lccccc}
    \toprule
    Level & Mesh file & Elements & $h_{\min}$ & $\Delta t$ & Status \\
    \midrule
    $h_0$
      & \texttt{periodic\_box\_structured\_2x1x1.vtk}
      & 768 & $4.4960{\times}10^{-2}$ & $3.5967{\times}10^{-2}$ & complete \\
    $h_0/2$
      & \texttt{periodic\_box\_structured\_nx16\_ny8\_nz8.vtk}
      & 6144 & $2.2480{\times}10^{-2}$ & $1.7984{\times}10^{-2}$ & complete \\
    $h_0/4$
      & \texttt{periodic\_box\_structured\_nx32\_ny16\_nz16.vtk}
      & 49152 & $1.1240{\times}10^{-2}$ & $8.9920{\times}10^{-3}$ & complete \\
    \bottomrule
  \end{tabular}
  }
\end{table}

\paragraph{Evidence.}
Figures~\ref{fig:conserved_p_invariants}--\ref{fig:conserved_p_errors} show the
$p$-refinement experiment.  The total electromagnetic energy is flat in time for
all polynomial orders over $T=100$, while the deviation from the analytical
energy decreases by several orders of magnitude as the order is increased.  The
same trend is observed for the linear and angular momenta.  The electric and
magnetic charge errors remain close to floating-point roundoff throughout the
run, which is consistent with the discrete periodic cancellation of the
divergence constraints.  The optical chirality is analytically zero for this
linearly polarized plane wave; the numerical value oscillates visibly at low
order and is strongly reduced by increasing $p$.

\begin{figure}[htbp]
  \centering
  \includegraphics[width=\textwidth]{../../output/1_invariants/plane_wave_strucutred/p_convergence/periodic_cuboid_cf_strcutred_invariants_panel.pdf}
  \caption{$p$-refinement study for the periodic travelling wave using centered
  fluxes.  The panels show the time evolution of total energy, optical
  chirality, linear momentum norm, angular momentum norm, and electric and
  magnetic charges.}
  \label{fig:conserved_p_invariants}
\end{figure}

\begin{figure}[htbp]
  \centering
  \includegraphics[width=\textwidth]{../../output/1_invariants/plane_wave_strucutred/p_convergence/periodic_cuboid_cf_strcutred_components_panel.pdf}
  \caption{$p$-refinement study for the component-wise linear and angular
  momenta.  The exact travelling wave has dominant $P_x$ and angular components
  $L_y$ and $L_z$; the remaining momentum components are expected to vanish.}
  \label{fig:conserved_p_components}
\end{figure}

\begin{figure}[htbp]
  \centering
  \includegraphics[width=\textwidth]{../../output/1_invariants/plane_wave_strucutred/p_convergence/periodic_cuboid_cf_structured_invariants_error_panel.pdf}
  \caption{$p$-refinement errors for the conserved-quantity diagnostics.  Scalar
  quantities are plotted as absolute errors, while momentum quantities are
  plotted as Euclidean errors of the integrated component vectors.}
  \label{fig:conserved_p_errors}
\end{figure}

Figures~\ref{fig:conserved_h_invariants}--\ref{fig:conserved_h_errors} show the
corresponding $h$-refinement experiment.  With $p=1$ fixed, reducing the mesh
size from $h_0$ to $h_0/4$ reduces the maximum-in-time energy error from about
$3.3\times10^{-1}$ to about $2.5\times10^{-2}$.  Over the same refinement range,
the maximum Euclidean error in the integrated linear momentum decreases from
about $6.2\times10^{-1}$ to about $2.7\times10^{-2}$, and the corresponding
angular-momentum error decreases from about $4.5\times10^{-1}$ to about
$1.9\times10^{-2}$.  Electric-charge errors stay below $5\times10^{-13}$ and
magnetic-charge errors stay below $10^{-10}$ on the plotted complete levels.

\begin{figure}[htbp]
  \centering
  \includegraphics[width=\textwidth]{../../output/1_invariants/plane_wave_strucutred/h_convergence/periodic_cuboid_cf_h_convergence_invariants_panel.pdf}
  \caption{$h$-refinement study for the periodic travelling wave at fixed
  $p=1$.  The plotted levels correspond to $h_0$, $h_0/2$, and $h_0/4$.}
  \label{fig:conserved_h_invariants}
\end{figure}

\begin{figure}[htbp]
  \centering
  \includegraphics[width=\textwidth]{../../output/1_invariants/plane_wave_strucutred/h_convergence/periodic_cuboid_cf_h_convergence_momentum_components_panel.pdf}
  \caption{$h$-refinement study for the component-wise linear and angular
  momenta at fixed $p=1$.  Mesh refinement reduces the spurious transverse
  momentum components and stabilizes the expected nonzero components.}
  \label{fig:conserved_h_components}
\end{figure}

\begin{figure}[htbp]
  \centering
  \includegraphics[width=\textwidth]{../../output/1_invariants/plane_wave_strucutred/h_convergence/periodic_cuboid_cf_h_convergence_invariants_error_panel.pdf}
  \caption{$h$-refinement errors for the conserved-quantity diagnostics at fixed
  $p=1$.  Energy and momentum errors decrease under mesh refinement; charge
  errors remain at roundoff scale.}
  \label{fig:conserved_h_errors}
\end{figure}

\paragraph{Interpretation.}
These experiments separate exact discrete conservation from accuracy of a
particular diagnostic.  The energy is the quadratic Hamiltonian of the linear
lossless Maxwell system and is the quantity most directly tied to the
skew-symmetry of the Poisson-bracket semi-discretization and to the symplectic
partitioned time integrator.  The observed non-drifting energy histories are
therefore the primary structure-preservation result.

The remaining diagnostics are global physical quantities evaluated from the
computed polynomial fields.  They are not all enforced as algebraic invariants
of the discrete bracket with the same strength as the Hamiltonian.  Their errors
therefore measure the quality of the discrete electromagnetic field, the
periodic face pairing, the quadrature evaluation, and the phase accuracy of the
long-time wave propagation.  The systematic reduction of the momentum and
chirality errors under both $p$- and $h$-refinement shows that these quantities
are recovered consistently as the DG approximation improves.  The charge
diagnostics being near roundoff is a stronger test of the distributed periodic
topology: it indicates that the periodic face contributions cancel globally and
that no artificial boundary leakage is introduced by the partitioned mesh.

Overall, the conserved-quantity study supports the central numerical claim of
the method: in a fully three-dimensional distributed-memory setting, the
Poisson-bracket DG formulation combined with ESPRK time integration delivers
long-time non-drifting energy behavior, while the physically relevant momentum,
charge, and chirality diagnostics converge toward their analytical values under
resolution refinement.

\end{document}
